% !TeX root = ../main.tex
\chapter{可查询 XML 压缩方法及问题分析}

\section{分析对象与比较维度}

第二章回顾的结构感知压缩、可查询压缩和 XML 索引研究，分别回答了如何利用 XML 结构组织压缩数据、如何在压缩表示上处理查询，以及如何借助索引减少访问范围等问题[4]--[17]。本章在这些成果基础上分析本文的研究对象，不再逐项介绍已有方法，而是聚焦一个贯穿编码、查询和索引构建的过程问题：编码阶段已经识别出的结构和值信息，能否以一致表示留存并被后续处理复用；复用这些信息所需的辅助空间及处理成本又如何衡量。

为分析该问题，本文从四个维度考察方法链：第一，信息形成，即 XML 解析过程中识别了哪些路径、值和节点关系；第二，信息留存，即哪些识别结果进入压缩表示或辅助结构；第三，信息消费，即查询函数和索引键抽取是否使用同一语义标识；第四，正确性约束，即摘要和哈希产生候选后如何恢复精确判断。该分析维度将压缩率、查询时间和索引成本置于同一处理链中，而不把任何单一组件孤立地视为方法贡献。

表~\ref{tab:prior-method-analysis} 将第二章介绍的代表性工作映射到上述分析维度。表中“本文分析重点”是本文据其技术路线提出的比较视角，并非对原论文实验结果的重新评价。

\begin{table}[htbp]
\centering
\small
\caption{代表性 XML 方法与本文的分析维度}
\label{tab:prior-method-analysis}
\begin{tabular}{p{0.22\linewidth}p{0.34\linewidth}p{0.34\linewidth}}
\toprule
代表性方法 & 相关研究关注点 & 本文据此分析的环节 \\
\midrule
XMill、XSC[4][5] & 利用 XML 结构或事件组织压缩数据 & 解析阶段识别的结构信息如何进入压缩表示 \\
XGrind、XPRESS[6][7] & 在压缩 XML 表示上支持查询处理 & 查询需要消费哪些压缩信息及其访问粒度 \\
XQueC、XSAQCT[9][10] & 根据查询需求、随机访问调整数据组织 & 查询能力与压缩单位、额外组织成本的关系 \\
DataGuides、IndexFabric、XQzip、SXSI[11][14]--[16] & 以结构摘要、路径键或压缩索引支持导航与检索 & 路径信息在文档内定位和跨文档索引中的职责 \\
\bottomrule
\end{tabular}
\end{table}

\section{编码阶段信息复用的分析}

XML 编码通常先经过解析或结构识别，再执行具体压缩。设输入文档为 $D$，解析阶段从中提取路径集合 $P_D$、文档序内容记录 $R_D$ 和节点结构记录 $T_D$。若这些结果只作为编码时的临时数据，压缩后丢弃，则后续查询需要从原始文档重新解析，或先解压再提取查询所需信息。若为查询另外建立派生结构，则需要重复识别路径和值，并承担独立存储及一致性维护成本。

另一种组织方式是令编码结果成为多个阶段共享的中间表示。对字符串路径 $p$ 和文本值 $v$ 分别定义映射 $\phi_p(p)$ 与 $\phi_v(v)$，将路径和值转换为标识；压缩内容流、路径辅助信息及索引键提取均消费这些标识。其直接作用不是保证每个步骤都免于扫描，而是避免对同一语义对象重复执行字符串识别和独立编码。特别是，GIN 建索引仍须读取压缩值并提取键；信息复用减少的是重新解析 XML 和重新定义 token 语义的需要，不能表述为零成本索引构建。

从这一分析可得出信息复用需要满足的三个条件：其一，编码端和查询端的标识映射具有一致的字典作用域；其二，压缩流保留查询及解码所需的顺序和结构关系；其三，派生的路径索引、Bloom 摘要和关系级倒排项具有明确的职责边界。若只保留叶路径和值，内部容器及节点边界可能不足以恢复；若只保留结构摘要，则又不能直接完成精确内容判断。因此，本文方法同时分析内容序列与结构序列，而不是把 XML 压平为无序键值集合。

\section{查询候选生成与精确验证分析}

查询处理可以分解为候选生成和精确判定两个阶段。令目标查询为 $q$，其真实匹配文档集合为 $A(q)$，候选生成阶段返回集合 $C(q)$。正确的过滤结构应满足 $A(q)\subseteq C(q)$，即可以多返回候选，但不能漏掉真实匹配；精确验证再从 $C(q)$ 中筛出 $A(q)$。这一关系适用于 Bloom Filter 及基于哈希的索引候选，但二者的使用范围不同。

在单文档内部，路径存在判断可由路径标识及文档级路径索引定位；路径取值还需要访问与该路径关联的文档序内容记录。Bloom Filter 针对内容值提供快速否定：若摘要判断不存在，可跳过该文档的精确内容扫描；若判断可能存在，则必须扫描内容标识或原文进行精确比较。由于 Bloom Filter 允许假阳性，阳性结果不是命中结论。在跨文档层面，GIN 根据内容键定位候选行；当长值以哈希键表示时，哈希碰撞可能引入额外候选，执行器仍需 recheck。

由此可以区分三种容易混淆的成本和作用：路径索引回答单个压缩文档中是否存在目标路径；文档级 Bloom 摘要减少单文档内容扫描；GIN 倒排结构减少关系中需要访问的行。它们可以协作，却不能互相替代。查询速度提升取决于候选缩减比例、精确复核成本、缓存与数据库执行器开销；仅由某一过滤结构的理论性质不能推出端到端延迟必然下降。

\section{压缩、查询与索引代价的耦合关系}

压缩与查询共享信息会形成空间—处理的耦合。设原始 XML 逻辑列值总空间为 $S_X$，压缩列值为 $S_C$，辅助表示为 $S_A$，关系级索引为 $S_I$，共享字典分摊空间为 $S_D$，则在某一明确统计口径下，完整持久化成本可写为
\[
S_{\mathrm{total}}=S_C+S_A+S_I+S_D.
\]
实际数据库 relation 统计可能已包含部分列值或 TOAST 开销，因此比较时必须按同一物理口径避免重复相加。该式强调：压缩列值变小并不自动推出含索引系统总空间变小；共享字典、GIN 和其他关系开销都应在同一边界内核算。

查询成本可按处理阶段表示为
\[
T_q=T_{\mathrm{lookup}}+T_{\mathrm{candidate}}+T_{\mathrm{scan}}+T_{\mathrm{verify}}+T_{\mathrm{executor}},
\]
其中各项分别表示查询标识映射、候选生成、压缩表示扫描、精确复核和数据库执行器开销。复用编码标识的目标是减少重复解析与字符串识别成本；Bloom 或 GIN 可能减少后续扫描量，但它们自身也有维护、访问和复核成本。因而，方法的实际效果应通过查询测量而非单项复杂度推断。

数据特征也影响压缩收益。设总叶记录数为 $N$，不同路径数为 $U_p$，不同短文本数为 $U_v$，结构模板数为 $U_s$。当文档数量增加而 $U_p$、$U_v$ 的增长较慢时，重复路径和值能被字典复用的次数增加；当单文档内容更密集时，固定头部、路径元数据及结构边界成本可由更多记录摊薄；当文档结构相似时，路径和节点关系的重复程度提高，结构标识更容易复用。因此，在字典和辅助信息策略相同的条件下，这些性质为压缩比随数据规模、密度和结构相似性提高而改善提供了机制解释。该判断是基于编码成本的趋势分析，仍需通过多组数据测量验证，不能替代实验结果。

\section{本文研究问题、适用范围与设计要求}

上述分析将本文研究问题归纳为：如何在 XML 结构感知压缩过程中复用产生的路径和值标识，使压缩表示、XPath 导向查询处理及索引键提取共享一致的信息基础，同时保留解码和精确验证所需的结构信息，并对由此产生的空间与查询代价进行可比评价？PostgreSQL 是本文实现和验证方法可行性的数据库平台；研究对象是压缩与索引优化方法，平台实现用于观察方法在数据库类型中的落地能力，而不是将研究贡献限定为 PostgreSQL 某种系统设计。

本文方法和实验遵循以下要求：一是使用结构感知表示压缩 XML，避免把层次化数据仅视为普通文本；二是令编码阶段产生的路径和值标识供压缩、查询及索引键抽取复用；三是明确路径判断、路径取值和内容查询的语义，并保留文档顺序及结构信息；四是将 Bloom 或哈希候选与精确结果判定分开；五是分别报告压缩列值、索引、relation、共享字典等空间口径，以及查询和写入测量的条件；六是把实验结论限定在当前原型、数据集和配置，不据此宣称对所有可查询 XML 压缩方法或数据库平台具有普遍优势。

\section{本章小结}

本章从作者的研究问题出发分析了编码信息复用的必要条件、候选生成与精确验证的关系，以及压缩、查询和索引成本之间的耦合。分析表明，方法核心在于以一致的编码标识连接压缩表示、查询处理和索引键抽取，同时用结构信息支持解码、用精确复核约束候选结果。第四章据此给出压缩表示及编码过程，第五章说明基于该表示的查询处理方法。
